\documentclass[12pt]{article}
\usepackage{amsmath,amssymb,amsfonts,amsthm}
\usepackage[margin=1in]{geometry}

\begin{document}

\begin{center}
{\Large\bfseries Chapter 4, Problem 18}\\[6pt]
Byron \& Fuller, \textit{Mathematics of Classical and Quantum Physics}
\end{center}

\bigskip

\textbf{Problem.} Given any linear transformation $A$ on a finite-dimensional vector space, prove that in a given coordinate system, the adjoint matrix always exists.

[Hint: Modify the reasoning of Theorem 4.7.]

\bigskip
\textbf{Solution.}

Let $V$ be an $n$-dimensional inner-product space with an orthonormal basis $\{e_1, e_2, \ldots, e_n\}$. Let $A$ be a linear transformation on $V$.

We need to show there exists a linear transformation $A^\dagger$ such that $(A^\dagger x, y) = (x, Ay)$ for all $x, y \in V$.

In the orthonormal basis, $A$ is represented by the matrix with entries $A_{ij} = (e_i, Ae_j)$. Define $A^\dagger$ by the matrix with entries:
\[
(A^\dagger)_{ij} = \overline{A_{ji}}.
\]
That is, $(A^\dagger)_{ij} = \overline{(e_j, Ae_i)} = (Ae_i, e_j)$.

We verify the defining property. For arbitrary vectors $x = \sum_i x_i e_i$ and $y = \sum_j y_j e_j$:
\begin{align*}
(x, Ay) &= \left(\sum_i x_i e_i,\; A\sum_j y_j e_j\right) = \sum_{i,j} \bar{x}_i y_j\,(e_i, Ae_j) = \sum_{i,j} \bar{x}_i y_j A_{ij}.
\end{align*}

Now compute $(A^\dagger x, y)$:
\begin{align*}
(A^\dagger x, y) &= \left(\sum_k (A^\dagger)_{ki} x_i\, e_k,\; \sum_j y_j e_j\right) = \sum_{k,i,j} \overline{(A^\dagger)_{ki} x_i}\, y_j\, (e_k, e_j) \\
&= \sum_{i,j} \overline{(A^\dagger)_{ji}}\, \bar{x}_i\, y_j = \sum_{i,j} \overline{\overline{A_{ij}}}\, \bar{x}_i\, y_j = \sum_{i,j} A_{ij}\, \bar{x}_i\, y_j.
\end{align*}

This equals $(x, Ay)$, confirming that $A^\dagger$ with matrix elements $(A^\dagger)_{ij} = \overline{A_{ji}}$ satisfies the required property.

Since the construction is explicit (conjugate transpose of the matrix), the adjoint always exists in any given coordinate system for a finite-dimensional space. \qed

\end{document}
